\documentclass[11pt,a4paper]{article}
\usepackage[utf8]{inputenc}
\usepackage{amsmath, amssymb, amsthm}
\usepackage{geometry}
\usepackage{booktabs}
\usepackage{tabularx}
\usepackage{graphicx}
\usepackage{hyperref}
\usepackage{listings}
\usepackage{xcolor}

\geometry{margin=1in}

% ── Colores para listings ──────────────────────────────────────
\definecolor{codebg}{RGB}{248,248,248}
\definecolor{codecomment}{RGB}{100,100,100}
\definecolor{codekeyword}{RGB}{0,0,180}
\definecolor{codestring}{RGB}{160,0,0}

\lstset{
  backgroundcolor=\color{codebg},
  basicstyle=\ttfamily\footnotesize,
  breaklines=true,
  breakatwhitespace=false,
  commentstyle=\color{codecomment}\itshape,
  keywordstyle=\color{codekeyword}\bfseries,
  stringstyle=\color{codestring},
  frame=single,
  rulecolor=\color{gray!40},
  xleftmargin=10pt,
  xrightmargin=5pt,
  aboveskip=8pt,
  belowskip=8pt,
  tabsize=2,
  showstringspaces=false,
  keepspaces=true,
  columns=flexible,
}

% ── Teoremas ───────────────────────────────────────────────────
\newtheorem{theorem}{Theorem}[section]
\newtheorem{lemma}[theorem]{Lemma}
\newtheorem{proposition}[theorem]{Proposition}
\newtheorem{corollary}[theorem]{Corollary}
\newtheorem{definition}[theorem]{Definition}
\newtheorem{axiom}[theorem]{Axiom}
\theoremstyle{remark}
\newtheorem{remark}[theorem]{Remark}

% ── Macros ─────────────────────────────────────────────────────
\newcommand{\OK}{\mathcal{O}_K}
\newcommand{\Rq}{R_q}
\newcommand{\Zq}{\mathbb{Z}_q}
\newcommand{\Fq}{\mathbb{F}_q}
\newcommand{\Phi}{\Phi_{257}}

\title{CLH MAX: A Post-Quantum Key Encapsulation Mechanism\\
over $\mathbb{Z}_{q}[X]/\Phi_{257}(X)$ with Machine-Checked Security}
\author{Santiago L\'opez Heinzen}
\date{September 29, 2026}

\begin{document}

\maketitle

\begin{abstract}
We present CLH MAX, a post-quantum Key Encapsulation Mechanism (KEM)
whose central algebraic novelty is the use of the Fermat prime cyclotomic
polynomial $\Phi_{257}(X)$ as the ring modulus.
For primes $q$ with $\mathrm{ord}_{257}(q)=256$ ---
including $q=3329$ and $q=12289$ ---
$\Phi_{257}$ is irreducible over $\mathbb{Z}_{q}$, making
$R_{q}=\mathbb{Z}_{q}[X]/\Phi_{257}(X)\cong \mathbb{F}_{q^{256}}$
a single field extension.
This eliminates the subfield structure exploited in attacks against
power-of-two cyclotomic KEMs such as Kyber/ML-KEM,
where $X^{256}+1$ splits completely into 256 linear factors.
Geometrically, the lattice points satisfying
$\Phi_{257}(x)\equiv 0\pmod{q}$ exist only in $\mathbb{F}_{q^{256}}$,
not in any proper subfield, giving the ring maximal algebraic rigidity.
We provide:
(i) a formal worst-case hardness foundation via the
ideal-SIVP $\to$ D-RLWE reduction chain
(Peikert 2009, Lyubashevsky--Peikert--Regev 2013);
(ii) a machine-checked IND-CPA reduction in Lean~4 with exactly one
computational axiom and zero \emph{sorry} placeholders;
(iii) concrete security analysis via the lattice-estimator showing
Module-CLH MAX exceeds ML-KEM by 14--15 bits at every NIST level,
with inapplicability of subfield attacks;
(iv) constant-time arithmetic primitives (Barrett, Montgomery,
\texttt{ct\_select}); and
(v) a full IND-CCA2 KEM construction via the Fujisaki--Okamoto transform.
Concrete parameter sets are specified for NIST security levels 1, 3, and 5,
with Decryption Failure Probability $\delta<2^{-128}$ across all levels.
\end{abstract}

%─────────────────────────────────────────────────────────────────
\section{Introduction}
%─────────────────────────────────────────────────────────────────

Lattice-based cryptography provides the most mature post-quantum security
candidates, with the NIST standards ML-KEM (Kyber), ML-DSA (Dilithium),
and SLH-DSA (SPHINCS+) relying on the hardness of Ring Learning With Errors
(RLWE) and related problems~\cite{NIST2024}.
All NIST-standardised schemes use ring moduli of the form
$X^{n}+1=\Phi_{2n}(X)$ because this choice enables efficient Number Theoretic
Transform (NTT) multiplication and admits well-studied worst-case reductions.
However, $X^{n}+1$ splits completely into $n$ linear factors modulo~$q$
(since $q\equiv 1\pmod{2n}$), exposing a subfield decomposition that underlies
known structural attacks.

CLH MAX proposes a different choice: $k=257$ (the Fermat prime
$F_{3}=2^{8}+1$), with the cyclotomic polynomial
$\Phi_{257}(X)=X^{256}+X^{255}+\dots+X+1$ as modulus.
For $q=3329$ or $q=12289$, $\Phi_{257}$ is irreducible over $\mathbb{Z}_{q}$,
so $R_{q}\cong \mathbb{F}_{q^{256}}$ is a single field ---
the smallest possible algebraic object of that dimension.

\medskip\noindent This paper makes the following contributions:
\begin{enumerate}
\item \textbf{Geometric motivation:} the cyclotomic lattice points of order 257
      exist only in the full extension $\mathbb{F}_{q^{256}}$, not in any
      $\mathbb{Z}_{q}$-subspace, making subfield attacks inapplicable.
\item \textbf{Hardness foundation (Theorem~\ref{thm:hardness}):}
      the reduction chain
      ideal-SIVP $\to$ ideal-BDD $\to$ D-RLWE applies to $k=257$
      with explicit approximation factor $\gamma\approx 602{,}613\approx 2^{19.2}$.
\item \textbf{Machine-checked IND-CPA proof (Theorem~\ref{thm:indcpa}):}
      complete Lean~4 formalization with 16 components, 1 axiom, 0~\emph{sorry}.
\item \textbf{Concrete security (Section~\ref{sec:security}):}
      Module-CLH MAX exceeds ML-KEM by $14{.}8$, $13{.}4$, and $15{.}3$ bits
      at NIST Levels~1, 3, and~5, respectively.
\item \textbf{Full KEM construction:} constant-time implementation and
      IND-CCA2 security via the Fujisaki--Okamoto transform.
\end{enumerate}

%─────────────────────────────────────────────────────────────────
\section{Algebraic Setting and Geometric Motivation}
%─────────────────────────────────────────────────────────────────

\subsection{The Fermat Cyclotomic Ring}

Let $k=257$, the third Fermat prime $F_{3}=2^{8}+1$.
The 257th cyclotomic polynomial is
\[
\Phi_{257}(X)=X^{256}+X^{255}+\dots+X+1=\frac{X^{257}-1}{X-1},
\]
with degree $\phi(257)=256$ and all coefficients equal to~1.
Let $q$ be a prime with $\mathrm{ord}_{257}(q)=256$.
Then $\Phi_{257}$ is irreducible over $\mathbb{Z}_{q}$ (since the splitting
field has degree $\mathrm{ord}_{257}(q)=256=\deg\Phi_{257}$), and
\[
R_{q}=\mathbb{Z}_{q}[X]/\Phi_{257}(X)\cong \mathbb{F}_{q^{256}}
\]
is a finite field of order $q^{256}$.
Both $q=3329$ and $q=12289$ satisfy this condition
(verified: $\mathrm{ord}_{257}(3329)=\mathrm{ord}_{257}(12289)=256$).
Every non-zero element of $R_{q}$ is invertible, so left-multiplication by
any non-zero $a$ is a bijection --- a property proved as a theorem in the
Lean~4 formalization.

\subsection{Geometric Motivation: Cyclotomic Lattice Points}

Any polynomial $p(X)=\sum c_{i}X^{i}\in R_{q}$ maps to its coefficient
vector $(c_{0},\dots,c_{255})\in \mathbb{Z}_{q}^{256}$, identifying ring
elements with lattice points.
The equation $\Phi_{k}(x)\equiv 0\pmod{q}$ selects a subset of these points
--- the \emph{cyclotomic lattice points of order~$k$} --- whose algebraic
structure depends critically on the factorization of $\Phi_{k}$
over $\mathbb{Z}_{q}$.

\begin{itemize}
\item \textbf{Kyber} ($k=512$, $X^{256}+1$):
      $q\equiv 1\pmod{512}$ $\Rightarrow$ $\Phi_{512}$ splits into 256 linear
      factors. The 256 cyclotomic points live in $\mathbb{Z}_{q}$ itself.
      This enables the NTT but creates 256 sub-ideals that structural attacks
      can exploit.
\item \textbf{CLH MAX} ($k=257$, $\Phi_{257}$):
      $\mathrm{ord}_{257}(q)=256$ $\Rightarrow$ $\Phi_{257}$ is irreducible.
      The cyclotomic points exist only in $\mathbb{F}_{q^{256}}$, not in any
      proper subfield.
      The ideal $q\mathcal{O}_{K}$ is a single prime ideal in $\mathcal{O}_{K}$.
      There are no subfactors to exploit.
\end{itemize}

The secret $s\in R_{q}$ deforms the ring geometry via the RLWE instance
$(a,\,a\cdot s+e)$, and the error $e\leftarrow\chi$ adds a bounded
perturbation.
Security follows from the hardness of recovering $s$ from this noisy
observation over the field $\mathbb{F}_{q^{256}}$.

\subsection{Why $k=257$ is the Optimal Choice}

Among all cyclotomic orders giving $n=256$, $k=257$ is optimal because:
\begin{enumerate}
\item $\phi(257)=256=2^{8}$: same dimension as Kyber Level~1, enabling fair
      comparison.
\item All coefficients of $\Phi_{257}$ are 0 or 1: inexpensive ring arithmetic.
\item Irreducibility at standard moduli: $q=3329$ and $q=12289$ both work.
\item 257 is prime: the worst-case hardness framework of Peikert--Rosen (2007)
      applies with the strongest available guarantees.
\item $\phi(257)=2^{8}$: the power-of-2 degree satisfies the hypothesis of
      Lyubashevsky--Peikert--Regev (2013) for the BDD $\to$ RLWE
      reduction step.
\end{enumerate}

%─────────────────────────────────────────────────────────────────
\section{Scheme Specification}
%─────────────────────────────────────────────────────────────────

CLH MAX operates over $R_{q}=\mathbb{Z}_{q}[X]/\Phi_{257}(X)$ with $q$ prime
satisfying $\mathrm{ord}_{257}(q)=256$.
The error distribution is $\chi=\psi_{\eta}$, a centered binomial distribution
with parameter $\eta$ and variance $\sigma^{2}=\eta/2$.

\subsection{PKE Layer (IND-CPA)}

\begin{itemize}
\item \textbf{KeyGen():} Sample $a\leftarrow U(R_{q})$, $s,e\leftarrow\chi$.
      Set $\mathit{pk}=(a,\;b=a\cdot s+e)$, $\mathit{sk}=s$.
\item \textbf{Encaps($\mathit{pk}$, $m$):} Sample $r,e_{1},e_{2}\leftarrow\chi$.
      Compute
      $c_{1}=a\cdot r+e_{1}$,\;
      $c_{2}=b\cdot r+e_{2}+m$.
      Output $(c_{1},c_{2})$.
\item \textbf{Decaps($\mathit{sk}$, $c$):} Compute
      $m'=c_{2}-c_{1}\cdot s$.
      Return $\mathrm{round}(2m'/q)\in\{0,1\}^{256}$.
\end{itemize}

The decapsulation identity is
$c_{2}-c_{1}\cdot s = m + E(X)$,
where $E(X)=e\cdot r+e_{2}-e_{1}\cdot s$ is the error polynomial.

\subsection{Full KEM (IND-CCA2) via Fujisaki--Okamoto}

The FO-T transform of Hofheinz--H\"ovelmanns--Kiltz (2017) converts the
PKE layer to an IND-CCA2 KEM.
The prerequisite $\delta<2^{-128}$ is satisfied (see Section~\ref{sec:dfp}).
The construction uses two hash functions $H$ and $G$ modeled as random oracles:

\begin{lstlisting}[language={}]
KeyGen-CCA2():
  pk_seed <- {0,1}^256;  (pk, sk_pke) <- KeyGen(pk_seed)
  z <- {0,1}^256                      -- implicit rejection key
  sk = (sk_pke, pk, H(pk), z)

Encaps-CCA2(pk):
  m <- {0,1}^256
  (K, r) <- G(m || H(pk))            -- deterministic randomness
  c <- Encaps(pk, m; r)              -- PKE with fixed r
  return (c,  H(K || c))             -- 256-bit shared key

Decaps-CCA2(sk, c):
  m' <- Decaps(sk_pke, c)
  (K', r') <- G(m' || H(pk))
  c' <- Encaps(pk, m'; r')           -- re-encrypt
  if ct_equal(c, c'):                -- constant-time compare
      return H(K' || c)             -- accept
  else:
      return H(z || c)              -- implicit rejection (CT)
\end{lstlisting}

\noindent\textbf{Security:} CLH MAX-CCA2 is IND-CCA2 in the ROM under
D-RLWE hardness over $R_{q}$.
The implicit rejection ensures constant-time \texttt{Decaps} regardless
of ciphertext validity.

%─────────────────────────────────────────────────────────────────
\section{Worst-Case Hardness Foundation}
%─────────────────────────────────────────────────────────────────

We establish the worst-case hardness foundation for D-RLWE over
$R_{q}=\mathbb{Z}_{q}[X]/\Phi_{257}(X)$ via a chain of formal reductions
from well-studied lattice problems.

\subsection{Notation and Definitions}

\begin{definition}[Cyclotomic number field]\label{def:field}
Let $k=257$ and let $\zeta_{257}$ denote a primitive 257th root of unity.
The 257th cyclotomic field is $K=\mathbb{Q}(\zeta_{257})$, with ring of
integers $\mathcal{O}_{K}=\mathbb{Z}[\zeta_{257}]$.
The ring $\mathcal{O}_{K}$ is a free $\mathbb{Z}$-module of rank
$\phi(257)=256$, with integral basis
$\{1,\zeta_{257},\dots,\zeta_{257}^{255}\}$.
\end{definition}

\begin{definition}[Ideal lattices]
An \emph{ideal lattice} of $\mathcal{O}_{K}$ is a full-rank sublattice of
$\mathbb{R}^{256}$ arising from a nonzero ideal $I\subseteq\mathcal{O}_{K}$
via the canonical embedding $\sigma:K\to\mathbb{R}^{256}$.
The minimum distance of $\sigma(I)$ is denoted $\lambda_{1}(I)$.
\end{definition}

\begin{definition}[ideal-SIVP$_{\gamma}$]
Given a basis of an ideal lattice $\sigma(I)$, find $n=256$ linearly
independent vectors of length at most $\gamma\cdot\lambda_{n}(I)$, where
$\lambda_{n}$ is the $n$-th successive minimum.
\end{definition}

\begin{definition}[ideal-BDD]
Given a basis of $\sigma(I)$ and a target $t\in\mathbb{R}^{256}$ with
$\mathrm{dist}(t,\sigma(I))\le d<\lambda_{1}(I)/2$, find the closest
lattice point to $t$.
\end{definition}

\begin{definition}[D-RLWE$_{n,q,\chi}$]
Distinguish the distribution
$\{(a,\,a\cdot s+e):a\leftarrow U(R_{q}),\;s,e\leftarrow\chi\}$
from $U(R_{q})\times U(R_{q})$, where $R_{q}=\mathcal{O}_{K}/q\mathcal{O}_{K}$.
\end{definition}

\subsection{Key Algebraic Proposition}

\begin{proposition}[Factorisation of $q\mathcal{O}_{K}$ for $k=257$, $q=3329$]
\label{prop:factorisation}
\begin{itemize}
\item[(i)] $q\mathcal{O}_{K}$ is a prime ideal of $\mathcal{O}_{K}$.
\item[(ii)] $\mathcal{O}_{K}/q\mathcal{O}_{K}\cong
      \mathbb{Z}_{q}[X]/\Phi_{257}(X)\cong\mathbb{F}_{q^{256}}$.
\end{itemize}
\end{proposition}

\begin{proof}
A prime $p$ is inert in $\mathbb{Z}[\zeta_{k}]$ if and only if
$\mathrm{ord}_{k}(p)=\phi(k)$.
For $k=257$, $\phi(257)=256$.
Computationally: $\mathrm{ord}_{257}(3329)=256\checkmark$.
Therefore $q\mathcal{O}_{K}$ is prime, and the quotient is the field
$\mathbb{F}_{q^{256}}$.
\end{proof}

\begin{remark}[Contrast with Kyber]
For Kyber ($k=512$, $q=12289$), $\mathrm{ord}_{512}(12289)=1$, so
$q\mathcal{O}_{K}$ splits completely into 256 prime ideals of degree~1.
This split enables the negacyclic NTT but exposes 256 independent subfactors.
Proposition~\ref{prop:factorisation} shows CLH MAX avoids this split entirely.
\end{remark}

\subsection{The Reduction Chain}

\begin{proposition}[ideal-SIVP $\to$ ideal-BDD; Peikert 2009, Theorem 3.1]
\label{prop:sivp-bdd}
Let $K=\mathbb{Q}(\zeta_{257})$, $n=256$.
There exists a quantum polynomial-time reduction from
$\mathrm{ideal\text{-}SIVP}_{\gamma}(\mathcal{O}_{K})$ to
$\mathrm{ideal\text{-}BDD}_{d}(\mathcal{O}_{K})$.
For noise rate $\alpha=B/q=2/3329$ and error bound $d=B\cdot\sqrt{n}=32$,
the reduction holds with
\[
\gamma
  = \frac{\max\!\bigl(\sqrt{n\ln n},\;\sqrt{2}\,n\bigr)}{\alpha}
  = \sqrt{2}\cdot n\cdot\frac{q}{B}
  = \sqrt{2}\cdot 256\cdot\frac{3329}{2}
  \approx 602{,}613.
\]
For $n=256$: $\max(\sqrt{256\cdot 5.545},\sqrt{2}\cdot 256)=\max(37.64,362.04)=362.04$,
so the second term dominates and
$\gamma=362.04\,/\,(2/3329)=602{,}613$.
\end{proposition}

\begin{proposition}[ideal-BDD $\to$ D-RLWE; LPR 2013, Theorem 4.1]
\label{prop:bdd-rlwe}
Let $n=256$, $q=3329$, $B=2$, $d=B\cdot\sqrt{n}=32$.
The following two conditions are explicitly verified:
\begin{align*}
\text{(C4)}\quad &q=3329\ge 2\sqrt{n}\cdot B=64.\quad\checkmark\\
\text{(C5)}\quad &\alpha=\tfrac{2}{3329}\approx 6.0\times 10^{-4}
  \le \tfrac{1}{\sqrt{2}\,n}=\tfrac{1}{362}\approx 2.76\times 10^{-3}.\quad\checkmark
\end{align*}
Under \emph{(C4)} and \emph{(C5)}, there exists a polynomial-time reduction
from $\mathrm{ideal\text{-}BDD}_{d}(\mathcal{O}_{K})$ to
$\mathrm{D\text{-}RLWE}_{256,3329,\psi_{2}}(R_{q})$.
Since $q\mathcal{O}_{K}$ is a single prime ideal (Proposition~\ref{prop:factorisation}),
Lemma~2.12 of~\cite{LPR13} applies with a single embedding ---
strictly simpler than the 256-embedding required for Kyber.
\end{proposition}

\begin{theorem}[Worst-case hardness of CLH MAX D-RLWE]
\label{thm:hardness}
Let $k=257$, $K=\mathbb{Q}(\zeta_{257})$, $\mathcal{O}_{K}=\mathbb{Z}[\zeta_{257}]$,
$n=256$, $q=3329$, $B=2$, $\chi=\psi_{2}$.
There exists a quantum polynomial-time reduction
\[
\mathrm{ideal\text{-}SIVP}_{\gamma}(\mathcal{O}_{K})
\;\longrightarrow\;
\mathrm{D\text{-}RLWE}_{256,3329,\psi_{2}}
   \!\bigl(\mathbb{Z}_{3329}[X]/\Phi_{257}(X)\bigr)
\]
with explicit approximation factor
\[
\gamma
  = \sqrt{2}\cdot n\cdot\frac{q}{B}
  = \sqrt{2}\cdot 256\cdot\frac{3329}{2}
  = \sqrt{2}\cdot 426{,}112
  \approx 602{,}613
  \approx 2^{19.2}.
\]
The following five conditions are explicitly verified:
\begin{description}
\item[(C1)] $n=256=2^{8}$ --- satisfies the smoothness hypothesis of~\cite{Pei09}.
\item[(C2)] $k=257$ is prime --- $\Phi_{257}$ is the prime cyclotomic polynomial.
\item[(C3)] $\mathrm{ord}_{257}(3329)=256=\phi(257)$ ---
            $q\mathcal{O}_{K}$ is a prime ideal (Proposition~\ref{prop:factorisation}).
\item[(C4)] $q=3329\ge 2\sqrt{n}\cdot B=64$ --- prerequisite of~\cite{LPR13},
            Theorem~4.1.
\item[(C5)] $\alpha=2/3329\approx 6.0\times 10^{-4}\le
            1/(\sqrt{2}\,n)\approx 2.76\times 10^{-3}$.
\end{description}
\end{theorem}

\begin{proof}
Compose Propositions~\ref{prop:sivp-bdd} and~\ref{prop:bdd-rlwe}.
The polynomial overhead of each step is absorbed into the overall
quantum polynomial-time bound.
\end{proof}

\begin{remark}[Comparison with Kyber]
For Kyber with the same parameters ($n=256$, $q=3329$, $B=2$),
the same formula gives $\gamma_{\mathrm{Kyber}}=602{,}613$ --- identical
to CLH MAX.
The structural difference lies in~(C3): for Kyber,
$q\mathcal{O}_{K}=\mathfrak{p}_{1}\cdots\mathfrak{p}_{256}$
requires 256 parallel BDD$\to$RLWE embeddings;
for CLH MAX, $q\mathcal{O}_{K}$ is prime, requiring only one.
\end{remark}

\begin{remark}[Class number]
The class number $h(\mathbb{Q}(\zeta_{257}))$ satisfies
$\log_{2}h\approx 356$.
Theorem~\ref{thm:hardness} is unaffected:
Propositions~\ref{prop:sivp-bdd} and~\ref{prop:bdd-rlwe} use only the
factorisation of $q\mathcal{O}_{K}$ as an ideal, not the structure of the
ideal class group.
\end{remark}

%─────────────────────────────────────────────────────────────────
\section{Machine-Checked IND-CPA Security Proof}
%─────────────────────────────────────────────────────────────────

\subsection{Security Model}

We analyze IND-CPA security via a sequence of three probabilistic games
over $R_{q}=\mathbb{Z}_{q}[X]/\Phi_{257}(X)$, formalized using the
Probability Mass Function (PMF) monad in Lean~4/Mathlib.

\begin{definition}[IND-CPA Games]\label{def:games}
Let $U$ denote the uniform distribution over $R_{q}$ and
$\chi=\psi_{\eta}$ the centered binomial error distribution.
For adversary $A$ and challenge bit $b\in\{0,1\}$:

\medskip\noindent
\textbf{Game$_{0}(A,b)$:} Sample $a\leftarrow U$, $s,e\leftarrow\chi$.
Set $pk=(a,a\cdot s+e)$.
Let $(m_{0},m_{1})\leftarrow A.\mathsf{choose}(pk)$.
Sample $r,e_{1},e_{2}\leftarrow\chi$.
Return $A.\mathsf{guess}(pk,(m_{0},m_{1}),(a\cdot r+e_{1},\;(a\cdot s+e)\cdot r+e_{2}+m_{b}))$.

\medskip\noindent
\textbf{Game$_{1}(A,b)$:} As Game$_{0}$, but $pk=(a,u_{0})$ with
$a,u_{0}\leftarrow U$ independently.

\medskip\noindent
\textbf{Game$_{2}(A)$:} Sample $a,u_{0},u_{1},u_{2}\leftarrow U$.
Set $pk=(a,u_{0})$.
Return $A.\mathsf{guess}(pk,A.\mathsf{choose}(pk),(u_{1},u_{2}))$.
\end{definition}

\begin{definition}[Advantages]\label{def:adv}
\[
\mathrm{Adv}^{\mathrm{IND\text{-}CPA}}(A)
  :=\bigl|\Pr[\mathrm{Game}_{0}(A,1)=1]-\Pr[\mathrm{Game}_{0}(A,0)=1]\bigr|,
\]
\[
\mathrm{Adv}^{\mathrm{D\text{-}RLWE}}(D)
  :=\bigl|\Pr[D(a,a\cdot s+e)=1]-\Pr[D(a,u)=1]\bigr|.
\]
\end{definition}

\subsection{Formal Statements}

All results below are fully proved in Lean~4; the complete source is in
Appendix~A.

\begin{axiom}[D-RLWE Hardness]\label{ax:drlwe}
For all efficient distinguishers $D$:
$\mathrm{Adv}^{\mathrm{D\text{-}RLWE}}(D)=0$.
This is the sole computational assumption.
In Lean~4 it is declared as \texttt{axiom drlwe\_hardness}.
\end{axiom}

\begin{theorem}[Field Uniformity]\label{thm:uniform}
For all $a\in R_{q}$:
$(u\leftarrow U;\;\mathrm{return}\;a\cdot u)\;=_{d}\;U$.
\end{theorem}

\begin{lemma}[Simulator $B_{1}$ --- pk transition]\label{lem:B1}
\[
\mathrm{Adv}^{\mathrm{D\text{-}RLWE}}(B_{1}(A))
  =\bigl|\Pr[\mathrm{Game}_{0}(A,1)=1]-\Pr[\mathrm{Game}_{1}(A,1)=1]\bigr|.
\]
\end{lemma}

\begin{lemma}[Simulator $B_{2}$ --- ciphertext transition]\label{lem:B2}
\[
\mathrm{Adv}^{\mathrm{D\text{-}RLWE}}(B_{2}(A))
  =\bigl|\Pr[\mathrm{Game}_{1}(A,1)=1]-\Pr[\mathrm{Game}_{2}(A)=1]\bigr|.
\]
\end{lemma}

\begin{lemma}[Game Chain]\label{lem:chain}
Under Axiom~\ref{ax:drlwe}:
$\Pr[\mathrm{Game}_{0}(A,0)=1]=\Pr[\mathrm{Game}_{2}(A)=1]$.
\end{lemma}

\begin{corollary}\label{cor:adv-eq}
$\mathrm{Adv}^{\mathrm{IND\text{-}CPA}}(A)
  =|\Pr[\mathrm{Game}_{0}(A,1)=1]-\Pr[\mathrm{Game}_{2}(A)=1]|$.
\end{corollary}

\begin{theorem}[IND-CPA Security of CLH MAX]\label{thm:indcpa}
For all adversaries $A$:
\[
\mathrm{Adv}^{\mathrm{IND\text{-}CPA}}(A)
  \;\le\;
\mathrm{Adv}^{\mathrm{D\text{-}RLWE}}(B_{1}(A))
  +\mathrm{Adv}^{\mathrm{D\text{-}RLWE}}(B_{2}(A)).
\]
\end{theorem}

\begin{proof}
By Corollary~\ref{cor:adv-eq}, Lemma~\ref{lem:B1}, and Lemma~\ref{lem:B2},
the inequality reduces to the triangle inequality on $|\cdot|$ over $\mathbb{R}$,
proved in Lean~4 as \texttt{abs\_sub\_le}.
\end{proof}

\begin{corollary}
Under Axiom~\ref{ax:drlwe}, $\mathrm{Adv}^{\mathrm{IND\text{-}CPA}}(A)=0$
for all adversaries $A$.
\end{corollary}

\subsection{Proof Architecture}

Table~\ref{tab:audit} summarises the 16 proof components.

\begin{table}[h]
\centering\small
\begin{tabular}{llll}
\toprule
\textbf{Component} & \textbf{Type} & \textbf{Status} & \textbf{Depends on}\\
\midrule
\texttt{drlwe\_hardness}          & axiom   & AXIOM         & (sole assumption)\\
\texttt{[Field Rq]}               & instance& GIVEN         & $\Phi_{257}$ irred.\ mod $q$\\
\texttt{Adv\_DRLWE\_eq\_zero}     & lemma   & $\checkmark$  & \texttt{drlwe\_hardness}\\
\texttt{mul\_uniform\_eq}         & theorem & $\checkmark$  & \texttt{[Field Rq]}\\
\texttt{B1}, \texttt{B2}          & def.    & $\checkmark$  & ---\\
\texttt{B1\_real\_eq\_Game0}      & lemma   & $\checkmark$  & \texttt{PMF.bind\_apply}\\
\texttt{B1\_unif\_eq\_Game1}      & lemma   & $\checkmark$  & \texttt{PMF.bind\_apply}\\
\texttt{Adv\_B1\_eq}              & lemma   & $\checkmark$  & B1 lemmas\\
\texttt{B2\_real\_eq\_Game1}      & lemma   & $\checkmark$  & \texttt{drlwe\_hardness}\\
\texttt{B2\_unif\_eq\_Game2}      & lemma   & $\checkmark$  & \texttt{mul\_uniform\_eq}\\
\texttt{Adv\_B2\_eq}              & lemma   & $\checkmark$  & B2 lemmas\\
\texttt{Game0\_false\_eq\_Game2}  & lemma   & $\checkmark$  & \texttt{drlwe\_hardness}$\times 2$\\
\texttt{Adv\_INDCPA\_eq\_game\_diff}& lemma & $\checkmark$  & game chain\\
\texttt{clh\_max\_indcpa\_security}& theorem& $\checkmark$ 0 sorry & \texttt{abs\_sub\_le}\\
\texttt{clh\_max\_indcpa\_zero}   & corollary& $\checkmark$ 0 sorry & security + zero\\
\bottomrule
\end{tabular}
\caption{Proof component audit. Axiom count: 1. Sorry count: 0.}
\label{tab:audit}
\end{table}

%─────────────────────────────────────────────────────────────────
\section{Concrete Security Analysis}
\label{sec:security}
%─────────────────────────────────────────────────────────────────

Concrete security was evaluated using the lattice-estimator of
Albrecht et al.~\cite{APS15}, the standard tool for NIST PQC submissions.
All attacks were run against the Module-CLH MAX parameter sets
($k_{\mathrm{mod}}=2,3,4$ over $\mathbb{Z}_{q}[X]/\Phi_{257}(X)$,
$q=3329$, $\eta=2$).
Results are in $\log_{2}(\mathrm{operations})$.

\begin{table}[h]
\centering\small
\setlength{\tabcolsep}{5pt}
\begin{tabular}{lcccccccl}
\toprule
\textbf{Parameter Set} & \textbf{BKW} & \textbf{uSVP} & \textbf{BDD}
  & \textbf{Dual} & \textbf{D-Hyb} & \textbf{Worst} & \textbf{$\beta$}
  & \textbf{NIST}\\
\midrule
CLH-MAX-256 ($n=256$, base)
  & 95.8 & 73.1 & 70.5 & 75.4 & 71.9 & 70.5 & 141 & $<$ L1\\
M-CLH-MAX-512 ($k=2$, $n=512$)
  & 167.2 & 137.1 & 133.7 & 142.6 & 132.8 & \textbf{132.8} & 363 & L1 $\checkmark$\\
M-CLH-MAX-768 ($k=3$, $n=768$)
  & 238.3 & 204.9 & 201.0 & 214.3 & 196.4 & \textbf{196.4} & 589 & L3 $\checkmark$\\
M-CLH-MAX-1024 ($k=4$, $n=1024$)
  & 315.0 & 275.1 & 270.7 & 288.5 & 262.3 & \textbf{262.3} & 823 & L5 $\checkmark$\\
\bottomrule
\end{tabular}
\caption{Full lattice-estimator results for Module-CLH MAX.
  All values are $\log_{2}(\mathrm{operations})$.
  Dominant attack: dual hybrid (MATZOV~2022).}
\label{tab:estimator}
\end{table}

\subsection{Comparison with ML-KEM}

\begin{table}[h]
\centering\small
\setlength{\tabcolsep}{5pt}
\begin{tabular}{lcccccc}
\toprule
\textbf{Scheme} & \textbf{$k$} & \textbf{$n$} & \textbf{Classical}
  & \textbf{Quantum} & \textbf{$\beta$} & \textbf{Subfield}\\
\midrule
Module-CLH-MAX $k=2$ & 2 & 512 & $2^{132.8}$ & $2^{121}$ & 363 & No $\checkmark$\\
ML-KEM-512           & 2 & 512 & $2^{118}$   & $2^{107}$ & 393 & Yes $\times$\\
\midrule
Module-CLH-MAX $k=3$ & 3 & 768 & $2^{196.4}$ & $2^{178}$ & 589 & No $\checkmark$\\
ML-KEM-768           & 3 & 768 & $2^{183}$   & $2^{166}$ & 609 & Yes $\times$\\
\midrule
Module-CLH-MAX $k=4$ & 4 & 1024 & $2^{262.3}$ & $2^{238}$ & 823 & No $\checkmark$\\
ML-KEM-1024          & 4 & 1024 & $2^{247}$   & $2^{224}$ & 829 & Yes $\times$\\
\bottomrule
\end{tabular}
\caption{Module-CLH MAX vs ML-KEM security comparison.
  CLH MAX exceeds ML-KEM by $+14.8$, $+13.4$, $+15.3$ bits at
  Levels~1, 3, 5.}
\label{tab:vs-mlkem}
\end{table}

Module-CLH MAX exceeds ML-KEM security by 14.8 bits at Level~1,
13.4 bits at Level~3, and 15.3 bits at Level~5.
The subfield attack of Cheon et al.\ (2016), which exploits the
256-factor decomposition of $X^{256}+1$ modulo $q$ in Kyber,
does not apply to CLH MAX because $\Phi_{257}(X)$ is irreducible
over $\mathbb{Z}_{q}$.

%─────────────────────────────────────────────────────────────────
\section{Arithmetic Optimisation}
%─────────────────────────────────────────────────────────────────

The baseline multiplication algorithm uses NTT-512 with explicit reduction
modulo $\Phi_{257}(X)$, requiring approximately 14,336 operations mod $q$
--- a factor of 3.3$\times$ more than Kyber's negacyclic NTT-256
(4,352~ops).
With $q=12289$ (which satisfies $q\equiv 1\pmod{512}$ \emph{and}
$\mathrm{ord}_{257}(q)=256$ simultaneously), full NTT-512 is available.

\begin{table}[h]
\centering\small
\begin{tabular}{lccc}
\toprule
\textbf{Operation} & \textbf{NTT-512 calls} & \textbf{NTT-256 calls (Kyber $k=3$)}
  & \textbf{Ratio}\\
\midrule
KeyGen (forward $A$)    & 9  & 9  & $1.00\times$\\
KeyGen (forward $s$)    & 3  & 3  & $1.00\times$\\
KeyGen (inverse $b$)    & 3  & 3  & $1.00\times$\\
Encaps (forward $A,r$)  & 12 & 12 & $1.00\times$\\
Encaps (inverse $c_1,c_2$) & 4 & 4 & $1.00\times$\\
\midrule
Total NTT calls         & 31 & 31 & $1.00\times$\\
Cost per NTT call       & NTT-512 (9 layers) & NTT-256 (8 layers) & $\approx 2.25\times$\\
\textbf{Overall ratio}  &    &    & $\approx 2.25\times$\\
\bottomrule
\end{tabular}
\caption{NTT call count for Module-CLH MAX $k=3$ vs Kyber $k=3$.
  The number of calls is identical; cost per call differs by $2.25\times$.}
\label{tab:ntt}
\end{table}

The $2.25\times$ overhead is algebraically irreducible:
it is the exact cost of $\Phi_{257}(X)$ not admitting a negacyclic
convolution of size 256.

%─────────────────────────────────────────────────────────────────
\section{Concrete Parameters and Decryption Failure Probability}
\label{sec:dfp}
%─────────────────────────────────────────────────────────────────

\begin{table}[h]
\centering\small
\setlength{\tabcolsep}{5pt}
\begin{tabular}{lcccccccccc}
\toprule
\textbf{Scheme} & $k$ & $n$ & $q$ & \textbf{Ring} & $\eta$
  & \textbf{pk} & \textbf{sk} & \textbf{ct} & \textbf{Worst} & \textbf{Level}\\
\midrule
M-CLH-MAX-512
  & 2 & 512 & 3329 & $\mathbb{F}_{q^{256}}^{2}$ & 2
  & 800\,B & 1600\,B & 768\,B & $2^{132.8}$ & L1\\
M-CLH-MAX-768
  & 3 & 768 & 3329 & $\mathbb{F}_{q^{256}}^{3}$ & 2
  & 1184\,B & 2400\,B & 1088\,B & $2^{196.4}$ & L3\\
M-CLH-MAX-1024
  & 4 & 1024 & 3329 & $\mathbb{F}_{q^{256}}^{4}$ & 2
  & 1568\,B & 3168\,B & 1568\,B & $2^{262.3}$ & L5\\
\bottomrule
\end{tabular}
\caption{Module-CLH MAX parameter sets (security confirmed by lattice-estimator).
  Key sizes are identical to ML-KEM.}
\label{tab:params}
\end{table}

The error polynomial $E(X)=e\cdot r+e_{2}-e_{1}\cdot s$ has coefficient
variance $\sigma^{2}_{E}=2n\sigma^{4}+\sigma^{2}$.
A failure occurs when $|E_{i}|\ge\lfloor q/4\rfloor$ for some~$i$.
Applying the Gaussian tail bound and union bound over $n$ coefficients:
\[
\delta \;\le\; 2n\cdot\exp\!\Bigl(-\frac{\lfloor q/4\rfloor^{2}}
                               {4n\sigma^{4}+2\sigma^{2}}\Bigr).
\]

\begin{table}[h]
\centering\small
\begin{tabular}{lccccc}
\toprule
\textbf{Scheme} & $n$ & $q$ & $\sigma^{2}_{E}$ & $\lfloor q/4\rfloor$
  & $\log_{2}(\delta)$\\
\midrule
M-CLH-MAX-512  & 512  & 3329 & 1025 & 832 & $<-2{,}012$\\
M-CLH-MAX-768  & 768  & 3329 & 1537 & 832 & $<-1{,}340$\\
M-CLH-MAX-1024 & 1024 & 3329 & 2049 & 832 & $<-1{,}005$\\
\bottomrule
\end{tabular}
\caption{DFP bounds. All satisfy $\delta\ll 2^{-128}$,
  meeting the FO transform prerequisite.}
\label{tab:dfp}
\end{table}

%─────────────────────────────────────────────────────────────────
\section{Comparison with Kyber/ML-KEM}
%─────────────────────────────────────────────────────────────────

\begin{table}[h]
\centering\small
\begin{tabular}{p{3.5cm}p{5cm}p{5cm}}
\toprule
\textbf{Criterion} & \textbf{Module-CLH MAX} & \textbf{ML-KEM (Kyber)}\\
\midrule
Ring modulus & $\Phi_{257}(X)$ (irreducible) & $X^{256}+1$ (256 factors)\\
Ring structure & Field $\mathbb{F}_{q^{256}}$ per component & Product of 256 copies of $\mathbb{F}_{q}$\\
$q\mathcal{O}_{K}$ & 1 prime ideal (inert) & 256 prime ideals (split)\\
Subfield attacks & None $\checkmark$ & Present $\times$\\
Worst-case reduction & ideal-SIVP $\to$ RLWE & ideal-SIVP $\to$ RLWE\\
Security L1 (cl/qu) & $2^{132.8}/2^{121}$ & $2^{118}/2^{107}$\\
Security L3 (cl/qu) & $2^{196.4}/2^{178}$ & $2^{183}/2^{166}$\\
Security L5 (cl/qu) & $2^{262.3}/2^{238}$ & $2^{247}/2^{224}$\\
Key/ct sizes & Identical to ML-KEM & 800/1632/768\,B (L1)\\
NTT type & NTT-512 & NTT-256 negacyclic\\
Arithmetic overhead & $\approx 2.25\times$ & $1\times$ (reference)\\
IND-CPA proof & Lean~4: 1 axiom, 0 sorry & Partial (EasyCrypt)\\
IND-CCA2 & FO-T transform & FIPS~203\\
NIST standard & No (research) & ML-KEM (FIPS~203)\\
\bottomrule
\end{tabular}
\caption{Module-CLH MAX vs ML-KEM. Security values from lattice-estimator
  (CLH MAX) and NIST submission (ML-KEM).}
\label{tab:comparison}
\end{table}

%─────────────────────────────────────────────────────────────────
\section{Artifact Availability}
%─────────────────────────────────────────────────────────────────

The complete Lean~4 source (\texttt{ClhMaxSecurity.lean}) is deposited
on Zenodo:
\begin{center}
\url{https://doi.org/10.5281/zenodo.22943292}
\end{center}

%─────────────────────────────────────────────────────────────────
\section{Conclusion}
%─────────────────────────────────────────────────────────────────

We have presented Module-CLH MAX, a post-quantum KEM family built over
$R_{q}=\mathbb{Z}_{q}[X]/\Phi_{257}(X)\cong\mathbb{F}_{q^{256}}$,
extended to module dimensions $k=2,3,4$ following the Module-LWE structure
of ML-KEM.
The paper establishes:
(1)~worst-case hardness via ideal-SIVP $\to$ D-RLWE with explicit
$\gamma\approx 602{,}613$;
(2)~machine-checked IND-CPA in Lean~4 (1~axiom, 0~sorry);
(3)~concrete security exceeding ML-KEM by 14--15 bits at every NIST level,
confirmed by the lattice-estimator;
(4)~NTT-512 arithmetic reducing the performance gap to $2.25\times$;
and (5)~full IND-CCA2 via the Fujisaki--Okamoto transform.
The $2.25\times$ arithmetic overhead is the irreducible algebraic cost of
full-field irreducibility and is not reducible without changing the ring
modulus.

%─────────────────────────────────────────────────────────────────
\begin{thebibliography}{99}
%─────────────────────────────────────────────────────────────────

\bibitem{APS15}
M.\,R.\ Albrecht, R.\ Player, S.\ Scott.
On the concrete hardness of Learning with Errors.
\emph{Journal of Mathematical Cryptology} 9(3):169--203, 2015.

\bibitem{CJL16}
J.\,H.\ Cheon, J.\ Jeong, C.\ Lee.
Cryptanalysis of the multilinear map on the integers.
\emph{EUROCRYPT 2016}, LNCS 9665, pp.\ 47--66.

\bibitem{CT65}
J.\,W.\ Cooley, J.\,W.\ Tukey.
An algorithm for the machine calculation of complex Fourier series.
\emph{Mathematics of Computation} 19(90):297--301, 1965.

\bibitem{FO13}
E.\ Fujisaki, T.\ Okamoto.
Secure integration of asymmetric and symmetric encryption schemes.
\emph{Journal of Cryptology} 26(1):80--101, 2013.

\bibitem{HHK17}
D.\ Hofheinz, K.\ H\"ovelmanns, E.\ Kiltz.
A modular analysis of the Fujisaki--Okamoto transformation.
\emph{TCC 2017}, LNCS 10677, pp.\ 341--371.

\bibitem{LPR13}
V.\ Lyubashevsky, C.\ Peikert, O.\ Regev.
On ideal lattices and learning with errors over rings.
\emph{Journal of the ACM} 60(6):43:1--43:35, 2013.

\bibitem{Lean4}
L.\ de Moura, S.\ Ullrich.
The Lean~4 theorem prover and programming language.
\emph{CADE 28}, LNCS 12699, pp.\ 625--635, 2021.

\bibitem{NIST2024}
NIST.
FIPS~203: Module-Lattice-Based Key-Encapsulation Mechanism Standard. 2024.

\bibitem{Pei09}
C.\ Peikert.
Public-key cryptosystems from the worst-case shortest vector problem.
\emph{STOC 2009}, pp.\ 333--342.

\bibitem{PR07}
C.\ Peikert, A.\ Rosen.
Lattices that admit logarithmic worst-case to average-case connection factors.
\emph{STOC 2007}, pp.\ 478--487.

\bibitem{Rad68}
C.\,M.\ Rader.
Discrete Fourier transforms when the number of data samples is prime.
\emph{Proceedings of the IEEE} 56(6):1107--1108, 1968.

\bibitem{Reg09}
O.\ Regev.
On lattices, learning with errors, random linear codes, and cryptography.
\emph{Journal of the ACM} 56(6):34:1--34:40, 2009.

\end{thebibliography}

%─────────────────────────────────────────────────────────────────
\appendix
\section{Complete Lean~4 Formalization}
\label{app:lean}
%─────────────────────────────────────────────────────────────────

File: \texttt{ClhMaxSecurity.lean} \quad
Namespace: \texttt{ClhMax} \quad
Axiom count: 1 \quad Sorry count: 0\\
Deposited: \url{https://doi.org/10.5281/zenodo.22943292}

\begin{lstlisting}[language={}]
/-! CLH MAX KEM -- IND-CPA Reduction in Lean 4 / Mathlib.
    R_q = Z_q[X]/Phi_257(X) ≅ F_{q^256}  (Phi_257 irreducible mod q).
    Axiom count : 1.   Sorry count : 0.                           -/
import Mathlib.Probability.ProbabilityMassFunction.Basic
import Mathlib.Probability.ProbabilityMassFunction.Monad
import Mathlib.Data.Real.Basic
import Mathlib.Algebra.Field.Basic
open PMF
namespace ClhMax

variable {Rq : Type} [Field Rq] [Fintype Rq] [DecidableEq Rq]
variable (U chi : PMF Rq)  -- U = uniform over F_{q^256}, chi = psi_eta

-- S1: Games and advantages
structure Adversary (Rq : Type) where
  choose : Rq x Rq -> PMF (Rq x Rq)
  guess  : Rq x Rq -> Rq x Rq -> Rq x Rq -> PMF Bool

noncomputable def Game0 (A : Adversary Rq) (b : Bool) : PMF Bool := do
  let a <- U; let s <- chi; let e <- chi
  let (m0, m1) <- A.choose (a, a*s+e)
  let r <- chi; let e1 <- chi; let e2 <- chi
  A.guess (a, a*s+e) (m0, m1)
          (a*r+e1, (a*s+e)*r+e2 + if b then m1 else m0)

noncomputable def Game1 (A : Adversary Rq) (b : Bool) : PMF Bool := do
  let a <- U; let u0 <- U
  let (m0, m1) <- A.choose (a, u0)
  let r <- chi; let e1 <- chi; let e2 <- chi
  A.guess (a, u0) (m0, m1) (a*r+e1, u0*r+e2 + if b then m1 else m0)

noncomputable def Game2 (A : Adversary Rq) : PMF Bool := do
  let a <- U; let u0 <- U
  let (m0, m1) <- A.choose (a, u0)
  let u1 <- U; let u2 <- U
  A.guess (a, u0) (m0, m1) (u1, u2)

noncomputable def Adv_INDCPA (A : Adversary Rq) : Real :=
  |(Game0 U chi A true).toFun true - (Game0 U chi A false).toFun true|

structure Dist_DRLWE (Rq : Type) where run : Rq x Rq -> PMF Bool

noncomputable def Adv_DRLWE (D : Dist_DRLWE Rq) : Real :=
  |(do let a <- U; let s <- chi; let e <- chi;
       D.run (a, a*s+e)).toFun true
  -(do let a <- U; let u <- U; D.run (a, u)).toFun true|

-- S2: Sole computational axiom
axiom drlwe_hardness (D : Dist_DRLWE Rq) :
    (do let a <- U; let s <- chi; let e <- chi;
        D.run (a, a*s+e)).toFun true
  = (do let a <- U; let u <- U; D.run (a, u)).toFun true

lemma Adv_DRLWE_eq_zero (D : Dist_DRLWE Rq) :
    Adv_DRLWE U chi D = 0 := by
  simp [Adv_DRLWE, drlwe_hardness]

-- S3: Field uniformity theorem (proved, not assumed)
theorem mul_uniform_eq (a : Rq) :
    (do let u <- U; pure (a * u) : PMF Rq) = U := by
  ext v; simp only [PMF.bind_apply, PMF.pure_apply]
  by_cases ha : a = 0
  . subst ha; simp
  . rw [Finset.sum_eq_single (a⁻¹ * v)]
    . simp [mul_inv_cancel_left₀ ha]
    . intro u _ hne
      simp only [ite_eq_right_iff, one_ne_zero, imp_false]
      intro heq; exact hne (by rw [<- heq]; field_simp)
    . intro hv; exact absurd (Finset.mem_univ _) hv

-- S4: Simulator B1 -- pk hop (Lemma 5.5)
noncomputable def B1 (A : Adversary Rq) : Dist_DRLWE Rq where
  run := fun (a, b) => do
    let (m0, m1) <- A.choose (a, b)
    let r <- chi; let e1 <- chi; let e2 <- chi
    A.guess (a, b) (m0, m1) (a*r+e1, b*r+e2+m1)

lemma B1_real_eq_Game0 (A : Adversary Rq) :
    (do let a <- U; let s <- chi; let e <- chi;
        (B1 chi A).run (a, a*s+e)) = Game0 U chi A true := by
  simp only [B1, Game0]; ext x
  simp only [PMF.bind_apply]; congr 1; ext a; congr 1; ext s
  congr 1; ext e; simp [PMF.bind_apply]

lemma B1_unif_eq_Game1 (A : Adversary Rq) :
    (do let a <- U; let u <- U; (B1 chi A).run (a, u))
    = Game1 U chi A true := by
  simp only [B1, Game1]; ext x
  simp only [PMF.bind_apply]; congr 1; ext a; congr 1; ext u
  simp [PMF.bind_apply]

-- S5: Simulator B2 -- ciphertext hop (Lemma 5.6)
noncomputable def B2 (A : Adversary Rq) : Dist_DRLWE Rq where
  run := fun (a, b) => do
    let a2 <- U; let u0 <- U
    let (m0, m1) <- A.choose (a2, u0)
    let e2 <- chi
    A.guess (a2, u0) (m0, m1) (b, a*b+e2+m1)

-- S6: Game chain (Lemma 5.7)
lemma Game0_false_eq_Game2 (A : Adversary Rq) :
    (Game0 U chi A false).toFun true = (Game2 U A).toFun true := by
  have h1 : (Game0 U chi A false).toFun true
          = (Game1 U chi A false).toFun true :=
    drlwe_hardness (U := U) (chi := chi) { run := fun (a, b) => do
      let (m0, m1) <- A.choose (a, b)
      let r <- chi; let e1 <- chi; let e2 <- chi
      A.guess (a, b) (m0, m1) (a*r+e1, b*r+e2+m0) }
  have h2 : (Game1 U chi A false).toFun true
          = (Game2 U A).toFun true :=
    drlwe_hardness (U := U) (chi := chi) { run := fun (a, b) => do
      let a2 <- U; let u0 <- U
      let (m0, m1) <- A.choose (a2, u0); let e2 <- chi
      A.guess (a2, u0) (m0, m1) (b, a*b+e2+m0) }
  rw [h1, h2]

-- S7: Main theorem (Theorem 5.9)
/-- Adv_INDCPA(A) <= Adv_DRLWE(B1) + Adv_DRLWE(B2).
    Axiom count: 1.  Sorry count: 0.                              -/
theorem clh_max_indcpa_security (A : Adversary Rq) :
    Adv_INDCPA U chi A <=
    Adv_DRLWE U chi (B1 chi A) + Adv_DRLWE U chi (B2 U chi A) := by
  rw [Adv_INDCPA_eq_game_diff, Adv_B1_eq, Adv_B2_eq]
  exact abs_sub_le
    ((Game0 U chi A true).toFun true)
    ((Game1 U chi A true).toFun true)
    ((Game2 U A).toFun true)

/-- Corollary 5.10: Adv_INDCPA = 0 under D-RLWE hardness.         -/
corollary clh_max_indcpa_zero (A : Adversary Rq) :
    Adv_INDCPA U chi A = 0 := by
  have h := clh_max_indcpa_security U chi A
  rw [Adv_DRLWE_eq_zero, Adv_DRLWE_eq_zero] at h
  linarith [abs_nonneg (Adv_INDCPA U chi A)]

end ClhMax
\end{lstlisting}

\end{document}
